% !TeX root = ../PIFB2.tex
\section{Introduction}\label{sec:introduction}

Wheeler's participatory program asks how registered distinctions acquire the significance of physical facts and how many situated observers come to describe one world~\cite{Wheeler1990}. A theory addressing that question needs more than a metaphor between information and matter. It needs laws for what an observer can register, how different descriptions can be compared, how collective descriptions form and what they predict. This monograph develops a probabilistic and geometric framework for those tasks, then examines what further work is required for physical geometry, persistence and time.

The starting point is local. An agent carries a probability law over a domain of an underlying continuous base. Where domains overlap, agents interact through selective links. Different points of the overlap carry different values of the same section fields and of the associated attention network. The base is not a graph of agents and is not presumed flat. Its ultimate physical interpretation remains open: the theory describes agents' records and models of underlying reality rather than presupposing that their descriptions exhaust it.

Two distinctions organize the construction. First, a coordinate identification between cognitive frames is not a communication channel. Sending may expose only a quotient of the sender's informational space, while receiving places that signal into a different subspace. Matrix-valued links, including rank-deficient links, represent this selective interaction. Invertible flat frame identifications are a special case. Second, the variational free energy of a normalized joint law is distinct from a composite objective assembled from pairwise divergences. Both can be studied, but they imply different probabilistic and dynamical claims.

We take the normalized joint as the foundation. Its collective evidence bound accommodates correlated beliefs, source-selection attention, model learning and observations that can be shared without being counted twice. Belief and model descriptions can live within one associated statistical bundle: a law over state and model coordinates admits a conditional state belief and a model marginal. This removes unnecessary duplication of primitive bundles. It does not explain time-scale separation by definition. We therefore examine actual mechanisms for accumulated memory, projected dynamics and slowly changing predictions.

At larger scales, a normalized coarse channel pushes fine laws to parent laws. Conditional expectation supplies the upward and reverse operators; integration generates new interactions and normalizing potentials; prediction supplies later priors. These constructions give precise versions of cross-scale influence without postulating an additional shadow field. Their remaining input is the retained variable or channel. Dynamical formation must explain why a particular retention becomes coherent, persistent and predictive, and why it can support the next formation episode.

The companion mathematical theory establishes a concrete selection result: at sufficiently small fixed temperature, slow source-prior learning on a homogeneous four-agent cycle selects one matching from almost every Gaussian initial mean, whereas a frozen-prior control has a symmetric long-time pairing law~\cite{DennisGaugeVFE2026}. The supplied initial mean provides the asymmetry. The result is a step toward endogenous organization, not a proof of general recursive hierarchy. Here its assumptions are stated alongside the exact parent laws and the remaining memory and closure obligations.

The physical part separates positive statistical geometry from Lorentz geometry. Fisher length measures distinguishability along belief histories, and Fisher pullbacks measure sensitivity to changes over the base. They are positive semidefinite. A different exact construction uses Hermitian two-by-two tensors: determinant polarization gives Lorentz signature, and determinant-preserving complex frame changes cover the proper orthochronous Lorentz group. Selective complex messages can generate these tensors classically. The mathematical representation is exact; selecting the required channel, a nondegenerate soldering to the base, physical scale and causal dynamics remains a physical research problem.

Persistence also admits a constructive treatment. Eliminating hidden first-order modes produces memory and colored fluctuations. It does not generally create Newtonian inertia. Gaussian calculations separate stiffness, friction, precision and covariance response, while a metastable example shows how a long residence time can arise without a kinetic-mass postulate. These calculations replace an identification of a rock with a high-precision belief by explicit response and lifetime questions.

\begin{table}[ht]
\centering\small
\begin{tabularx}{\textwidth}{@{}p{0.22\textwidth}XX@{}}
\toprule
Participatory question & Mathematical construction & Remaining identification\\
\midrule
What is registered? & Normalized record channels; expected information gain. & A physical registration mechanism and its irreversibility.\\
How do perspectives meet? & Selective covariant links and common observables. & Empirical meaning and calibration of the channels.\\
How do collectives form? & Learned selection, exact retained laws and memory. & Recursively inherited formation under general interfaces.\\
How do scales act back? & Conditional inference, learned parameters and prediction. & A dynamical choice of retained variables and boundary laws.\\
How does physical geometry arise? & Fisher geometry; conditional determinant-cone Lorentz kinematics. & Soldering, scale, causal characteristics and field equations.\\
\bottomrule
\end{tabularx}
\caption{Wheeler's motivation and the corresponding research obligations. The third column is not implied by the second.}\label{tab:wheeler_correspondence}
\end{table}

\paragraph{Interpretive stance.}\label{sec:cognitive_first}
The motivating stance is participatory and relational: physical descriptions are made through interactions, and collective descriptions should be accountable to the records of their constituents. This stance does not require treating a prior as a belief about nothing, denying an underlying reality, or assigning consciousness to every mathematical section. Calling a boundary mechanism an agent changes no probability law unless a criterion of agency changes its predictions. The monograph therefore asks what agency can explain rather than using universal agency as an explanatory premise.

The philosophical ambition is retained: stable things might be understood through persistent patterns of constrained participation. The mathematical work establishes conditional mechanisms and identifies the quantities they preserve. Whether those mechanisms describe matter is a separate question to be answered through dynamics and observation.

\paragraph{Reading route.}
Part I develops the continuous-base setting, the collective variational principle, learning and memory, and exact coarse-graining at one base point. Part II examines information clocks, induced geometry, the Lorentz construction and physical interpretation. The appendices retain full Gaussian calculations and normalized hierarchical examples. They are specializations whose assumptions are stated explicitly, not an implementation specification.

